% =========================================================
% Metric Spaces — Curated Exercise Set (Foundational Only)
% Sources: Abbott, Bartle & Sherbert, Bruckner, Tao, Lebl
% =========================================================

\section*{Metric Spaces — Exercise Review Set}

These exercises develop the foundational structure of metric spaces.
They focus strictly on the metric itself, open balls, boundedness,
sequences, and Cauchy sequences. Topics such as compactness and
continuity are intentionally excluded.

\bigskip

% ---------------------------------------------------------
% Stub Macro
% ---------------------------------------------------------

\newcommand{\exstub}[4]{
\subsection*{Exercise: #1}

\textbf{Source:} #2

\textbf{Technique:} #3

\textbf{Task.} #4

\begin{remark*}[Work Area]
Write your proof or computation here.
\end{remark*}

\bigskip
}

% =========================================================
% ABBOTT
% =========================================================

\section*{Exercises from Abbott}

\exstub
{Abbott Metric 1}
{Abbott, \textit{Understanding Analysis}}
{metric verification}
{Verify that the discrete metric satisfies the metric axioms.}

\exstub
{Abbott Metric 2}
{Abbott}
{metric verification}
{Verify that $d(x,y)=|x-y|$ defines a metric on $\mathbb R$.}

\exstub
{Abbott Metric 3}
{Abbott}
{triangle inequality}
{Prove the triangle inequality for the standard metric on $\mathbb R$.}

\exstub
{Abbott Metric 4}
{Abbott}
{ball containment}
{Show that every open ball in a metric space is an open set.}

\exstub
{Abbott Metric 5}
{Abbott}
{distance computation}
{Compute the distance between two points in $\mathbb R^2$ under the Euclidean metric.}

\exstub
{Abbott Metric 6}
{Abbott}
{counterexample}
{Give an example of a function that fails the triangle inequality and
therefore is not a metric.}

\exstub
{Abbott Metric 7}
{Abbott}
{metric construction}
{Construct a bounded metric on $\mathbb N$.}

\exstub
{Abbott Metric 8}
{Abbott}
{diameter bound}
{Define the diameter of a subset of a metric space and compute the
diameter of an interval.}

% =========================================================
% BARTLE
% =========================================================

\section*{Exercises from Bartle \& Sherbert}

\exstub
{Bartle Metric 1}
{Bartle \& Sherbert}
{metric verification}
{Verify that the taxicab metric defines a metric on $\mathbb R^2$.}

\exstub
{Bartle Metric 2}
{Bartle \& Sherbert}
{distance computation}
{Compute distances in $\mathbb R^2$ under the maximum metric.}

\exstub
{Bartle Metric 3}
{Bartle \& Sherbert}
{triangle inequality}
{Prove the triangle inequality for the maximum metric.}

\exstub
{Bartle Metric 4}
{Bartle \& Sherbert}
{ball containment}
{Describe the open balls in the taxicab metric on $\mathbb R^2$.}

\exstub
{Bartle Metric 5}
{Bartle \& Sherbert}
{diameter bound}
{Show that a subset of a metric space is bounded iff its diameter is finite.}

\exstub
{Bartle Metric 6}
{Bartle \& Sherbert}
{metric construction}
{Construct a new metric from an existing metric using
$d'(x,y)=\min(d(x,y),1)$.}

\exstub
{Bartle Metric 7}
{Bartle \& Sherbert}
{counterexample}
{Give an example of a pseudometric that fails the identity axiom.}

\exstub
{Bartle Metric 8}
{Bartle \& Sherbert}
{ball containment}
{Show that if $y$ lies in an open ball centered at $x$ then there exists
a smaller ball around $y$ contained in the first ball.}

% =========================================================
% BRUCKNER
% =========================================================

\section*{Exercises from Bruckner}

\exstub
{Bruckner Metric 1}
{Bruckner, \textit{Real Analysis}}
{metric verification}
{Verify that $d(x,y)=|x-y|$ is a metric on $\mathbb R$.}

\exstub
{Bruckner Metric 2}
{Bruckner}
{metric construction}
{Define a bounded metric on $\mathbb R$ equivalent to the standard metric.}

\exstub
{Bruckner Metric 3}
{Bruckner}
{distance computation}
{Compute the diameter of a closed interval.}

\exstub
{Bruckner Metric 4}
{Bruckner}
{sequence convergence}
{Define convergence of a sequence in a metric space.}

\exstub
{Bruckner Metric 5}
{Bruckner}
{sequence convergence}
{Show that limits of sequences in metric spaces are unique.}

\exstub
{Bruckner Metric 6}
{Bruckner}
{sequence convergence}
{Show that every convergent sequence is bounded.}

\exstub
{Bruckner Metric 7}
{Bruckner}
{Cauchy criterion}
{Define a Cauchy sequence in a metric space.}

\exstub
{Bruckner Metric 8}
{Bruckner}
{triangle inequality}
{Use the triangle inequality to prove convergent sequences are Cauchy.}

% =========================================================
% TAO
% =========================================================

\section*{Exercises from Tao}

\exstub
{Tao Metric 1}
{Tao, \textit{Analysis I}}
{metric verification}
{Verify that the discrete metric satisfies the metric axioms.}

\exstub
{Tao Metric 2}
{Tao}
{distance computation}
{Compute distances under the Euclidean metric.}

\exstub
{Tao Metric 3}
{Tao}
{triangle inequality}
{Prove the triangle inequality for the Euclidean metric.}

\exstub
{Tao Metric 4}
{Tao}
{metric construction}
{Construct a bounded metric equivalent to the Euclidean metric.}

\exstub
{Tao Metric 5}
{Tao}
{sequence convergence}
{Define convergence of sequences in metric spaces.}

\exstub
{Tao Metric 6}
{Tao}
{sequence convergence}
{Show that convergent sequences are bounded.}

\exstub
{Tao Metric 7}
{Tao}
{Cauchy criterion}
{Define Cauchy sequences and prove convergent sequences are Cauchy.}

\exstub
{Tao Metric 8}
{Tao}
{counterexample}
{Give an example of a Cauchy sequence that does not converge
in $\mathbb Q$.}

% =========================================================
% LEBL
% =========================================================

\section*{Exercises from Lebl}

\exstub
{Lebl Metric 1}
{Lebl, \textit{Basic Analysis I}}
{metric verification}
{Verify that the Euclidean metric satisfies the metric axioms.}

\exstub
{Lebl Metric 2}
{Lebl}
{distance computation}
{Compute distances in $\mathbb R^n$.}

\exstub
{Lebl Metric 3}
{Lebl}
{ball containment}
{Describe open balls in $\mathbb R^n$.}

\exstub
{Lebl Metric 4}
{Lebl}
{sequence convergence}
{Define convergence of sequences in metric spaces.}

\exstub
{Lebl Metric 5}
{Lebl}
{Cauchy criterion}
{Define Cauchy sequences and prove basic properties.}

\exstub
{Lebl Metric 6}
{Lebl}
{sequence convergence}
{Show limits of sequences are unique in metric spaces.}

\exstub
{Lebl Metric 7}
{Lebl}
{diameter bound}
{Define and compute the diameter of subsets of $\mathbb R^n$.}

\exstub
{Lebl Metric 8}
{Lebl}
{metric construction}
{Construct a bounded metric equivalent to the standard metric.}
